% astro-geo-backend_tech_notes.tex
\documentclass[11pt,a4paper]{article}

\usepackage[a4paper,margin=2.3cm]{geometry}
\usepackage[T1]{fontenc}
\usepackage[utf8]{inputenc}
\usepackage{lmodern}
\usepackage{microtype}
\usepackage{hyperref}
\usepackage{enumitem}
\usepackage{xcolor}
\usepackage{booktabs}
\usepackage{listings}

\hypersetup{
  colorlinks=true,
  linkcolor=blue!60!black,
  urlcolor=blue!60!black,
  citecolor=blue!60!black
}

\lstdefinestyle{bash}{
  basicstyle=\ttfamily\small,
  columns=fullflexible,
  frame=single,
  breaklines=true
}
\lstdefinestyle{python}{
  basicstyle=\ttfamily\small,
  columns=fullflexible,
  frame=single,
  breaklines=true
}

\title{astro-geo-backend --- Technical Notes (Current Stage)}
\author{Alessandro Veltri}
\date{\today}

\begin{document}
\maketitle

\section{Scope and non-scope}
This repository is a small backend whose job is to convert:
\begin{enumerate}[label=\arabic*.]
  \item a \textbf{place query} (e.g.\ a city name),
  \item a \textbf{local civil date and time},
\end{enumerate}
into a single deterministic JSON payload containing:
\begin{itemize}
  \item WGS84 coordinates and an IANA timezone identifier,
  \item the corresponding UTC instant (with explicit DST ambiguity handling),
  \item minimal astronomical context: \textbf{IAU constellation at the zenith} and \textbf{Sun position/constellation}.
\end{itemize}

\textbf{Non-scope:} this project explicitly does \emph{not} implement any astrological interpretation, horoscope logic, house systems, or narrative mapping. It only produces clean, reproducible inputs for downstream modules.

\section{High-level pipeline}
The current backend is a three-stage pipeline:
\[
\texttt{GEO} \rightarrow \texttt{TIME} \rightarrow \texttt{ASTRO}
\]
or in words:
\begin{enumerate}[label=\textbf{\alph*)}]
  \item \textbf{GEO:} resolve a place query to $(\phi,\lambda)$ and a tzid,
  \item \textbf{TIME:} resolve (local date, local time, tzid) to a \emph{unique} UTC instant,
  \item \textbf{ASTRO:} compute zenith constellation and Sun context at that UTC instant.
\end{enumerate}

The orchestration is performed by a single CLI entry point, which prints \textbf{exactly one JSON object} to stdout.

\section{Repository structure (current)}
The relevant components are organized under \texttt{src/}:
\begin{verbatim}
src/
|-- geo/         (place -> coords -> tzid)
|-- civil_time/  (local -> UTC, DST-safe)
|-- astro/       (zenith constellation, Sun position/constellation)
`-- main.py      (GEO -> TIME -> ASTRO orchestration, JSON output)
\end{verbatim}

\section{Stage 1: GEO --- place query to coordinates (+ tzid)}
\subsection{Data source and constraints}
The geocoding stage uses OpenStreetMap Nominatim via HTTPS GET queries.
Because Nominatim is a shared public service, the module implements:
\begin{itemize}
  \item \textbf{caching} of resolved queries to a JSON file in \texttt{\$HOME/.cache/},
  \item \textbf{rate limiting} (minimum delay between remote requests),
  \item a \textbf{user-agent string} (required by Nominatim usage policy).
\end{itemize}

\subsection{Optional timezone resolution}
Timezone resolution from coordinates is performed only if the optional Python dependency
\texttt{timezonefinder} is installed. If it is not available:
\begin{itemize}
  \item the GEO stage returns coordinates but \textbf{tzid is missing},
  \item the main pipeline stops early with a clear error telling the user to install \texttt{timezonefinder}.
\end{itemize}

\subsection{Cache behavior}
A cache hit skips the network entirely. Additionally, old cache entries can be \emph{upgraded}:
if an entry has no timezone stored and \texttt{timezonefinder} is now available, tzid is computed and stored.

\subsection{Inputs and outputs}
\textbf{Input:} place query string.\\
\textbf{Output:} best match (currently \texttt{limit=1}) with:
\[
\{\texttt{display\_name},\ \texttt{lat},\ \texttt{lon},\ \texttt{tzid}\}.
\]

\section{Stage 2: TIME --- local civil time to unique UTC instant}
This stage solves a subtle problem: local civil times are not always one-to-one with UTC.

\subsection{DST fall-back ambiguity (two valid instants)}
During the autumn DST transition, a local wall-time can happen twice.
The time module detects this by attempting both PEP~495 folds and verifying round-trips.

Policy options:
\begin{itemize}
  \item \texttt{ambiguous='raise'}: error out (strict mode),
  \item \texttt{ambiguous='earliest'}: choose the earlier UTC instant,
  \item \texttt{ambiguous='latest'}: choose the later UTC instant.
\end{itemize}

\subsection{DST spring-forward gaps (non-existent local times)}
During the spring DST transition, a local wall-time can be invalid.
Policy options:
\begin{itemize}
  \item \texttt{nonexistent='raise'}: error out (strict mode),
  \item \texttt{nonexistent='shift\_forward'}: move forward in small steps until valid,
  \item \texttt{nonexistent='shift\_backward'}: move backward in small steps until valid.
\end{itemize}

\subsection{Why this matters}
Astronomical computations must be fed an \textbf{unambiguous instant}.
This module ensures the pipeline always reaches a unique UTC time, while also allowing the CLI to:
\begin{itemize}
  \item start strict (raise),
  \item fall back to a chosen policy,
  \item emit a \textbf{warning string} into the JSON payload when a policy is applied.
\end{itemize}

\section{Stage 3: ASTRO --- minimal astronomical context}
All astronomical computations are performed with \texttt{astropy}.

\subsection{Zenith constellation}
Given:
\[
\texttt{utc\_iso},\ \phi=\texttt{lat\_deg},\ \lambda=\texttt{lon\_deg},
\]
the code constructs an \texttt{AltAz} frame at that Earth location and time, defines the zenith
as $(\texttt{alt}=90^\circ,\ \texttt{az}=0^\circ)$ (azimuth is arbitrary at zenith but fixed for determinism),
transforms to ICRS, and queries the official IAU constellation boundaries.

Output can be either:
\begin{itemize}
  \item full constellation name, or
  \item 3-letter IAU abbreviation (if \texttt{--short-constellation} is passed).
\end{itemize}

\subsection{Sun position and constellation}
At the same UTC instant, the Sun is obtained via \texttt{astropy.coordinates.get\_sun}.
The code then returns:
\begin{itemize}
  \item Sun constellation (IAU boundaries),
  \item true ecliptic of date longitude/latitude,
  \item ICRS right ascension and declination.
\end{itemize}

\section{The main CLI: single JSON payload}
\subsection{Command-line interface}
The primary entry point is:
\begin{lstlisting}[style=bash]
python3 src/main.py --city "Cosenza, Italy" --date 1982-08-25 --time 12:00
\end{lstlisting}

\subsection{What it prints}
The program prints \textbf{one JSON object} with three top-level blocks:
\begin{itemize}
  \item \texttt{place}: \texttt{display\_name}, \texttt{lat}, \texttt{lon}, \texttt{tzid}
  \item \texttt{time}: \texttt{local}, \texttt{utc}, \texttt{warnings}
  \item \texttt{astro}: zenith constellation + Sun block
\end{itemize}

This contract (one JSON object to stdout) is intentional: it allows the CLI to be used in pipes,
scripts, and future services.

\section{Determinism and reproducibility}
\subsection{What is deterministic}
\begin{itemize}
  \item Given \textbf{fixed} \texttt{(lat,lon,tzid,local time)}, the TIME stage is deterministic.
  \item Given \textbf{fixed} \texttt{(lat,lon,utc)}, the ASTRO stage is deterministic.
  \item Given \textbf{cached} geocoding entries, GEO is deterministic (no network).
\end{itemize}

\subsection{What is not fully deterministic}
\begin{itemize}
  \item The GEO stage depends on an external service (Nominatim) \emph{unless cached}.
  \item Place-name ambiguity is resolved by ``best result'' with \texttt{limit=1}, which can change
        if Nominatim ranking changes or the query is underspecified.
\end{itemize}

\section{Strengths and flaws (current stage)}
\subsection{Strengths}
\begin{itemize}
  \item \textbf{Clean separation of concerns:} GEO, TIME, ASTRO are independent and testable.
  \item \textbf{DST-safe time resolution:} ambiguity and gaps are detected explicitly, with policies.
  \item \textbf{Single structured output:} one JSON object is easy to integrate downstream.
  \item \textbf{Graceful optional dependency:} timezone resolution is modular (\texttt{timezonefinder}).
  \item \textbf{IAU correctness:} constellation lookups use official IAU boundaries (via Astropy).
\end{itemize}

\subsection{Flaws / known limitations}
\begin{itemize}
  \item \textbf{Geocoding is external:} without cache, GEO depends on network + Nominatim availability.
  \item \textbf{Timezone resolution is approximate:} \texttt{timezonefinder} uses polygons; edge cases near borders
        or in complex administrative regions can fail or be wrong.
  \item \textbf{Underspecified place queries:} ``Parma'' vs ``Parma, Ohio'' shows that tzid can be missing or
        ambiguous depending on the match quality.
  \item \textbf{Single-result geocoding:} \texttt{limit=1} hides alternatives; this is fine for a backend but not ideal
        for interactive UX.
  \item \textbf{No uncertainty metadata:} the JSON does not currently carry a confidence score beyond Nominatim
        ``importance'' (and even that is not surfaced to the user by default).
  \item \textbf{Astro scope is minimal:} zenith constellation and Sun context only; no Moon, planets, or
        local sky visibility constraints.
\end{itemize}

\section{Practical next steps (suggested)}
If the goal is to turn this into a robust backend for a larger app, the immediate improvements that
have the highest payoff are:
\begin{enumerate}[label=\arabic*.]
  \item \textbf{Expose geocoding alternatives:} allow \texttt{--limit > 1} and return candidates
        (or add an interactive/selection mode in a separate CLI).
  \item \textbf{Persist last-request time in cache or state:} current orchestration resets \texttt{last\_t} per run,
        so rate limiting is only meaningful inside a single process run.
  \item \textbf{Add a validation/diagnostics mode:} e.g.\ include zenith RA/Dec, UTC offset, DST flag in JSON when
        \texttt{--debug} is passed.
  \item \textbf{Add unit tests:} especially for DST edge cases (Europe/Rome and at least one non-European zone).
  \item \textbf{Define a stable JSON schema:} even a simple \texttt{schema.json} + version field prevents pain later.
\end{enumerate}

\section{Appendix: example JSON (conceptual)}
\begin{lstlisting}[style=python]
{
  "place": { "display_name": "...", "lat": 0.0, "lon": 0.0, "tzid": "Europe/Rome" },
  "time":  { "local": "YYYY-MM-DDTHH:MM", "utc": "YYYY-MM-DDTHH:MMZ", "warnings": [] },
  "astro": {
    "zenith_constellation": "Monoceros",
    "sun": {
      "constellation": "Sagittarius",
      "ecliptic": { "lon_deg": 271.6, "lat_deg": -0.003 },
      "icrs": { "ra_deg": 309.6, "dec_deg": -19.44 }
    }
  }
}
\end{lstlisting}

\end{document}
